% zfs.tex — ZFS / OpenZFS, the overview: how a pool is built from vdevs and datasets, never
% overwriting + checksums + self-healing, snapshots and what they cost, send/receive, the write and
% read paths, and the choices that cannot be undone. Facts ranked by importance, as many as fit;
% the subtopics (pool layout, snapshots & replication, write path, ARC, root & boot, integrity,
% encryption & dedup) get folios of their own.
% Sources (checked 2026-10-11) — first checked 2026-10-06:
%   github.com/openzfs/zfs man/man7/zpoolconcepts.7 ("Data is dynamically distributed across all
%     top-level devices"; a vdev faults with "Insufficient replicas"; raidz "can sustain one, two, or
%     three failures"; dRAID "integrated distributed hot spares", "fixed stripe width … to allow
%     fully sequential resilvering"; special = "all metadata, the indirect blocks of user data, intent
%     log (in absence of separate log device), and deduplication tables", small blocks opt-in, full
%     special class spills to normal; "Log devices can be mirrored"; "Cache devices cannot be mirrored
%     or part of a raidz configuration", a cache read error "is reissued to the original storage pool
%     device"; hot spares replace a failed device automatically; checkpoint = "a pool-wide snapshot",
%     rewinding "will permanently discard it")
%   man/man7/zfsconcepts.7 (snapshots "created extremely quickly, and initially consume no additional
%     space", "As data within the active dataset changes, the snapshot consumes more data"; snapshots
%     "automatically mounted on demand"; clones "can only be created from a snapshot", promote reverses
%     the dependency; bookmarks reference "when a snapshot was created", no space)
%   man/man7/zfsprops.7 ("Properties are inherited from the parent unless overridden by the child";
%     a snapshot's used = "space that is referenced exclusively by this snapshot", freed if destroyed;
%     written = "written since the previous snapshot"; copies "in addition to any redundancy",
%     "stored on different disks, if possible", charged to the dataset; compression on = lz4 when the
%     feature is enabled; recordsize "affects only files created afterward"; sync=disabled "is very
%     dangerous"; encryption on = aes-256-gcm; dedup off)
%   man/man8/zpool-scrub.8 (scrub "examines all data … and verifies each block's checksum"; replicated
%     devices: "ZFS automatically repairs any damage"; resilver "only examines data that ZFS knows to
%     be out of date")
%   man/man8/zfs-send.8 (incremental -i; from a bookmark; -w raw "may then be received on an untrusted
%     machine", cannot be read "or alter[ed] without being detected"; -t resumes)
%   man/man8/zfs-receive.8 (-s keeps partial state, receive_resume_token; -c corrective receive,
%     "Metadata can not be healed")
%   man/man8/zpool-remove.8 (raidz/draid top-level vdevs block removal; ashift must match) ·
%     man/man8/zpool-attach.8 (raidz expansion: "old blocks retain their old data-to-parity ratio") ·
%     man/man7/zpool-features.7 (an unsupported active feature blocks read-write import, and read-only
%     import unless read-only compatible; compatibility limits zpool upgrade)
%   zfsprops.7 also: quota "a dataset and its descendants", refquota excludes "descendants,
%     including file systems and snapshots", reservation = "minimum amount of space guaranteed"
%   man/man4/zfs.4 (zfs_arc_max=0: the larger of RAM − 1 GiB and 5/8 of RAM; auto ashift) ·
%     module/zfs/txg.c (zfs_txg_timeout = 5) · OpenZFS docs, Workload Tuning (ashift 12 for 4K disks)
%   en.wikipedia.org/wiki/OpenZFS (CDDL vs GPL → out-of-tree module on Linux); measured on OpenZFS
%     2.3.9, Debian 13: the module is built per kernel (DKMS) and must be in the initramfs for a ZFS
%     root; <mountpoint>/.zfs/snapshot/<name>/ reads back a file with snapdir=hidden
% @labels: area=os kind=concept platform=general topic=storage,data discussion=141
% @tags: zfs, openzfs, copy-on-write, merkle-tree, raidz, draid, mirror, zil, slog, arc, l2arc, snapshots, send-receive, checksums, scrub, compression, ashift
\documentclass[8pt]{extarticle}
\usepackage{printup-sheet}
\usepackage{array}

% Linked text: the website links it, paper prints the text only (paper holds no links).
\newcommand\link[2]{\pdfonly{#1}\htmlonly{\href{#2}{#1}}}
\newcommand\ct[1]{\texttt{#1}}
\newcommand\hd[1]{\par\vspace{4pt}\noindent{\bfseries\color{sheetBlue}#1}\par\vspace{2pt}}

\tikzset{
  ttl/.style={font=\bfseries, text=sheetBlue, anchor=west},
  ds/.style={box, draw=sheetBrown, fill=sheetBrown!8, font=\small\ttfamily, inner sep=2pt},
  pool/.style={box, draw=sheetBlue, fill=sheetBlue!8, font=\small\bfseries, inner sep=2.5pt},
  vd/.style={box, draw=sheetGreen!60!black, fill=sheetGreen!12, font=\small\bfseries, inner sep=2pt},
  dk/.style={cell, minimum width=4.6mm, minimum height=4mm, fill=sheetGreen!8},
  kind/.style={box, draw=sheetOrange, fill=sheetOrange!10, font=\small\bfseries, inner sep=2pt,
               minimum width=19mm},
  say/.style={font=\small, align=left, anchor=west, inner sep=1pt},
  lbl/.style={font=\small, text=black!75, align=center, inner sep=1pt},
  blk/.style={box, font=\small, inner sep=1.5pt, minimum width=9mm, minimum height=4.5mm},
  old/.style={blk, draw=sheetGrey, fill=black!5, text=black!65},
  new/.style={blk, draw=sheetGreen!60!black, fill=sheetGreen!14},
  held/.style={blk, draw=sheetOrange, fill=sheetOrange!15},
  wrote/.style={blk, draw=sheetBlue, fill=sheetBlue!12},
  stage/.style={box, font=\small, inner sep=2pt, minimum height=5mm, text width=24mm, align=center},
}

\begin{document}
\pdfonly{\folioslug{os/zfs}}

\sheettitle{ZFS — pools, copy-on-write, checksums, snapshots}{os · folio}

\oneliner{ZFS is a volume manager and a filesystem in one. Disks form \textbf{vdevs}, vdevs form a
\textbf{pool}, and every \textbf{dataset} draws on the pool's free space. Nothing is overwritten in
place: each change goes to new blocks under a new root, so a crash leaves the old state or the new
one, a \textbf{snapshot} costs nothing to take, and a disk that returns wrong data is caught by a
\textbf{checksum} — and, with redundancy, repaired.}

\vspace{3pt}
\noindent\begin{tikzpicture}[sheet]
  % ───────── the pool
  \node[ttl] at (0,4.55) {How ZFS is built};
  \node[lbl, anchor=west] at (2.3,4.55) {· datasets inherit properties down the tree; a child may override};
  \node[ds] (h) at (1.1,3.8) {tank/home};
  \node[ds] (d) at (3.3,3.8) {tank/db};
  \node[ds] (v) at (5.75,3.8) {tank/vm (a zvol)};
  \node[pool, minimum width=88mm] (p) at (4.4,3.0) {pool \ct{tank} — every dataset draws on the same free space};
  \foreach \n in {h,d,v} \draw[flow] (\n) -- (\n |- p.north);
  \node[vd] (m) at (1.6,2.05) {mirror};
  \node[vd] (r) at (6.0,2.05) {raidz2};
  \draw[flow] (p.south -| m) -- (m); \draw[flow] (p.south -| r) -- (r);
  \node[lbl] at (3.8,2.32) {data is spread over\\every top-level vdev};
  \node[dk] at (1.3,1.3) {}; \node[dk] at (1.9,1.3) {};
  \foreach \i in {0,...,5} \node[dk] at (4.7+0.52*\i,1.3) {};
  \draw (m.south) -- (1.6,1.52); \draw (r.south) -- (6.0,1.52);
  \node[lbl] at (1.6,0.82) {survives all but one disk};
  \node[lbl] at (6.0,0.82) {survives any 2 of the 6};
  \node[lbl, text=sheetRed, anchor=west] at (0,0.22)
    {Redundancy lives \textbf{inside each vdev}: lose one vdev, lose the \textbf{whole pool}.};
  % ───────── the other kinds of vdev
  \draw[sheetGrey!40] (9.25,4.75) -- (9.25,0.0);
  \node[ttl] at (9.45,4.55) {Other kinds of vdev};
  \foreach \y/\k/\t in {
      3.95/{raidz1, raidz3}/{1 or 3 parity disks: survive 1 or 3 failed disks},
      3.23/dRAID/{spare space spread over all disks, a fixed stripe: a sequential, faster rebuild},
      2.51/special/{metadata and the dedup table (small blocks opt-in). \textcolor{sheetRed}{Lose it, lose the pool}: mirror it},
      1.79/log (SLOG)/{the intent log: only \emph{sync} writes get faster. Can be mirrored},
      1.07/cache (L2ARC)/{more read cache, no redundancy needed: a failed read goes to the pool},
      0.35/spare/{takes over a failed disk on its own}} {
    \node[kind] at (10.45,\y) {\k};
    \node[say, text width=50mm] at (11.5,\y) {\t};
  }
\end{tikzpicture}

\begin{multicols}{3}
\raggedright\setstretch{1.0}

\hd{Never overwrite}
\noindent\begin{tikzpicture}[sheet]
  \node[old] (r0) at (0.7,2.0) {root $n$};
  \node[old] (i0) at (0.7,1.25) {pointers};
  \node[old] (b) at (0.3,0.45) {B};
  \node[new] (r1) at (3.75,2.0) {root $n{+}1$};
  \node[new] (i1) at (3.75,1.25) {pointers$'$};
  \node[new] (b1) at (4.15,0.45) {B$'$};
  \node[blk, draw=black!60, fill=white] (a) at (2.2,0.45) {A};
  \draw[->, sheetGrey] (r0) -- (i0); \draw[->, sheetGrey] (i0) -- (b); \draw[->, sheetGrey] (i0) -- (a);
  \draw[->, thick, sheetGreen!60!black] (r1) -- (i1); \draw[->, thick, sheetGreen!60!black] (i1) -- (b1);
  \draw[->, thick, sheetGreen!60!black] (i1) -- (a);
  \node[lbl] at (2.2,-0.15) {A is shared, not copied};
\end{tikzpicture}\par
Changing B writes B$'$ and new pointers up to a new root, \textbf{written last}. A crash leaves the
old tree or the new one: no half-written state, no \ct{fsck}.

Every pointer holds the \textbf{checksum} of the block it points to, all the way up to the root. A
disk that returns wrong data, or the right data from the wrong place, fails the check.

On a mirror, raidz or dRAID the bad copy is \textbf{rewritten from a good one}. A \textbf{scrub}
reads everything to find damage before you need the data; a \textbf{resilver} rebuilds only what
is out of date. \ct{copies=2} keeps a second copy even on a single disk.

\columnbreak

\hd{Snapshots}
\noindent\begin{tikzpicture}[sheet]
  \node[lbl, anchor=east] at (0.55,1.6) {\ct{@mon}};
  \node[new] at (1.05,1.6) {A}; \node[held] at (1.95,1.6) {B}; \node[held] at (2.85,1.6) {C};
  \node[lbl, anchor=east] at (0.55,0.85) {now};
  \node[new] at (1.05,0.85) {A}; \node[wrote] at (1.95,0.85) {B$'$};
  \node[lbl, text=black!50] at (2.85,0.85) {deleted};
  \node[lbl, anchor=west, align=left, text=sheetGreen!45!black] at (3.4,1.9) {A: shared};
  \node[lbl, anchor=west, align=left, text=sheetOrange!80!black] at (3.4,1.45) {B, C: held only\\by \ct{@mon}};
  \node[lbl, anchor=west, align=left, text=sheetBlue] at (3.4,0.75) {B$'$: written\\since \ct{@mon}};
\end{tikzpicture}\par
A snapshot is \textbf{the old root kept alive}: instant, nothing copied, no space at first. It grows
as the live data changes away from it.

Deleting C frees nothing while \ct{@mon} holds it. Destroying \ct{@mon} frees only its \ct{used}:
the blocks nothing else holds (B and C). \ct{written} is what changed since the last snapshot.
A \ct{quota} caps a dataset with its children and snapshots, a \ct{refquota} the dataset alone; a
\ct{reservation} guarantees it space.

Any snapshot reads back at \ct{.zfs/snapshot/<name>/}, even when hidden: one file back, no rollback.
A \textbf{clone} is a writable dataset made from a snapshot, sharing its blocks.

\columnbreak

\hd{Send and receive}
\noindent\begin{tikzpicture}[sheet]
  \node[lbl, font=\small\ttfamily] at (0.75,2.05) {tank/db};
  \node[old] (s1) at (0.75,1.4) {@1}; \node[new] (s2) at (0.75,0.65) {@2};
  \node[lbl, font=\small\ttfamily] at (4.0,2.05) {backup/db};
  \node[old] (t1) at (4.0,1.4) {@1}; \node[new, dashed] (t2) at (4.0,0.65) {@2};
  \draw[hot] (s2.east) -- node[lbl, above, text=sheetOrange!80!black]{only blocks} node[lbl, below, text=sheetOrange!80!black]{born after @1} (t2.west);
  \node[lbl] at (2.4,-0.05) {over a pipe or \ct{ssh}};
\end{tikzpicture}\par
\ct{send} turns a snapshot, or \textbf{only the difference between two}, into a stream;
\ct{receive} replays it into another pool. The cost follows what changed, not the size. An
incremental send can also start from a \textbf{bookmark}, which marks when a snapshot was taken and
costs nothing, so the old snapshot itself can go.

A \textbf{raw} send stays encrypted: the backup host stores it without being able to read it, or
to change it unnoticed. An interrupted receive kept with \ct{-s} resumes from its token. A
corrective receive repairs damaged data from a healthy stream.

\textcolor{sheetRed}{\textbf{A snapshot is not a backup}}: it lives on the same disks. A backup is a
snapshot sent somewhere else.

\end{multicols}

\begin{multicols}{3}
\raggedright\setstretch{1.0}

\hd{Writing}
\noindent\begin{tikzpicture}[sheet]
  \node[stage, draw=sheetGrey, fill=white] (w) at (1.3,2.05) {\ct{write()}, \ct{fsync()}};
  \node[stage, draw=sheetBlue, fill=sheetBlue!8] (t) at (1.3,1.2) {RAM: transaction group};
  \node[stage, draw=sheetGreen!60!black, fill=sheetGreen!12] (pl) at (1.3,0.35) {pool, every $\le$5 s};
  \node[box, draw=sheetOrange, fill=sheetOrange!10, font=\small, inner sep=2pt, text width=13mm, align=center] (z) at (4.05,1.2) {ZIL\\(on SLOG)};
  \draw[flow] (w) -- (t); \draw[flow] (t) -- (pl);
  \draw[hot] (w.east) -| node[lbl, pos=0.25, above]{sync} (z.north);
\end{tikzpicture}\par
Writes gather in RAM and land together. A sync write (\ct{fsync}, a database commit) also goes to
the \textbf{intent log}, which is read only after a crash. A \textbf{SLOG} makes sync writes faster
and nothing else; \ct{sync=disabled} acknowledges writes it has not stored.

\ct{recordsize} is the largest block of a dataset and applies to files created afterwards.
Compression is on by default (lz4): fewer bytes to write and read.

\columnbreak

\hd{Reading}
\noindent\begin{tikzpicture}[sheet]
  \node[stage, draw=sheetBlue, fill=sheetBlue!8] (arc) at (1.3,2.05) {ARC, in RAM};
  \node[stage, draw=sheetOrange, fill=sheetOrange!10] (l2) at (1.3,1.2) {L2ARC, on SSD};
  \node[stage, draw=sheetGreen!60!black, fill=sheetGreen!12] (pl) at (1.3,0.35) {the pool's disks};
  \draw[flow] (arc) -- node[lbl, right]{miss} (l2); \draw[flow] (l2) -- node[lbl, right]{miss} (pl);
\end{tikzpicture}\par
The \textbf{ARC} caches what is read most often and most recently. By default it may grow to the
larger of RAM $-$ 1 GiB and 5/8 of RAM: on a ZFS box, ``used'' memory is mostly cache.

An \textbf{L2ARC} extends it onto an SSD, but its index lives in RAM. Dedup is off by default for
the same reason: its table must stay in RAM to be fast. Encryption is chosen per dataset.

\columnbreak

\hd{Choices you cannot undo}
\begin{itemize}
  \trap{\ct{ashift}, the sector size, is fixed per vdev when it is created: \ct{12} for disks with 4K
        sectors.}
  \trap{A raidz or dRAID vdev can never be removed from the pool.}
  \trap{A \textbf{checkpoint} rewinds the whole pool, configuration included — and rewinding
        discards it.}
  \trap{A single disk added beside a raidz becomes the pool's weakest vdev.}
  \trap{Expanding a raidz with a disk keeps old blocks at their old data:parity ratio.}
  \trap{\ct{zpool upgrade} can make the pool unreadable to an older ZFS; the \ct{compatibility}
        property limits it.}
  \trap{On Linux, ZFS is a separate kernel module, rebuilt for every kernel; a ZFS root needs it in
        the initramfs.}
\end{itemize}

\begin{htmlonly}
\section{Further reading}
\begin{itemize}
  \item \href{https://openzfs.github.io/openzfs-docs/}{OpenZFS documentation} — the project's own guides, from getting started to performance tuning
  \item \href{https://github.com/openzfs/zfs/blob/master/man/man7/zpoolconcepts.7}{zpoolconcepts(7)} — every kind of vdev, what each survives, and how a pool fails
  \item \href{https://github.com/openzfs/zfs/blob/master/man/man7/zfsconcepts.7}{zfsconcepts(7)} — datasets, snapshots, clones and bookmarks, in the manual's words
\end{itemize}
\end{htmlonly}

\end{multicols}

\noindent{\footnotesize\color{sheetGrey}\textit{Related:} linux-filesystems-storage · git-in-depth
(a Merkle DAG too) · sqlite-grdb-internals (WAL vs ZIL) · os-fundamentals}

\end{document}
